\documentclass[10pt,a4paper]{article}

% ── Modern font setup ────────────────────────────────────────
\usepackage{fontspec}
\usepackage[ngerman]{babel}
\setmainfont{LiberationSans}[
  Extension      = .ttf,
  UprightFont    = *-Regular,
]
\renewcommand{\familydefault}{\sfdefault}

\usepackage[margin=1.5cm,top=1.8cm,bottom=1.5cm]{geometry}
\usepackage{fancyhdr}
\usepackage{lastpage}
\usepackage{tabularx}
\usepackage{tabularray}
\usepackage{microtype}
\usepackage[hidelinks]{hyperref}
% \usepackage{stackengine}
\usepackage{array}
\usepackage{enumitem}
\usepackage{graphicx}
\usepackage{amssymb}
\usepackage[table]{xcolor}

% \setlength{\fboxsep}{4pt}
\setlength{\parindent}{0pt}
\setlength{\parskip}{0pt}
\renewcommand{\arraystretch}{1.05}
\setlist[itemize]{leftmargin=*,itemsep=1pt,topsep=2pt,parsep=0pt}

\newcommand{\cb}{\(\square\)}
\newcommand{\cbx}{\(\boxtimes\)}

% Form section heading: bold text, NOT boxed (original uses plain bold text)
\newcommand{\formsection}[1]{%
    \par\vspace{6pt}%
    \noindent\textbf{#1}%
    \par\vspace{6pt}%
}

% Subsection inside a bordered box
\newcommand{\formsubsection}[1]{\par\vspace{2pt}\textbf{#1}\par\vspace{2pt}}

% Form field underline (fixed width)
\newcommand{\formline}[1]{\underline{\hspace{#1}}}
% Form field underline (fills remaining space in column)
% \newcommand{\formlinefill}{%
%   \TextField[
%     name=\the\pdfelapsedtime, % Unique name for each field
%     width=\linewidth,         % Fills the cell width
%     borderstyle=U,            % Underline style
%     height=10pt               % Adjust to match text height
%   ]{}%
% }
\newcommand{\formlinefill}{\leavevmode\xleaders\hbox{\rule[-.5pt]{1pt}{0.4pt}}\hfill\kern0pt}

\pagestyle{fancy}
\fancyhf{}
\lhead{}
\rhead{}
\rfoot{{\small Seite \thepage}}
\cfoot{}
\renewcommand{\headrulewidth}{0pt}
\renewcommand{\footrulewidth}{0pt}

\begin{document}
\begin{Form}

% === HEADER ===
\noindent
\begin{minipage}[t]{0.5\textwidth}
    \vspace{0pt}
    {\Large\textbf{Erhebungsbogen}}
\end{minipage}%
\hfill
\begin{minipage}[t]{0.45\textwidth}
  \raggedleft
  \vspace{0pt}
  {\LARGE\textbf{Enzkreis}}
\end{minipage}

\vspace{16pt}

% === AKTENZEICHEN BOX ===
\begin{tblr}{
  width = \textwidth,
  colspec = { |X| },
  row{1} = {abovesep = 20pt, belowsep = 0pt, valign = m},
  row{2} = {abovesep = 0pt},
}
  \hline
  Aktenzeichen: \hspace{0.2cm}
  \TextField[name=aktenzeichen, width=5cm, borderstyle=U]{}
  \\[1pt]
  \SetCell{font=\small}
  (falls bekannt / wird ansonsten von Ihrem Jobcenter ausgefüllt)
  \\ \hline
\end{tblr}

\vspace{6pt}

% === TEIL 1 HEADING (no box, just bold text with en-dash) ===
\formsection{Teil 1 \textendash{} Persönliche Daten}

% === PERSOENLICHE INFORMATIONEN (bordered box containing header + fields) ===
\begin{tblr}{
  width = \textwidth,
  colspec = { |lX[wd=5cm]X| }
}
  \hline
  \SetCell[c=3]{l, font=\bfseries}
    Persönliche Informationen
  \\[11pt]
  Vorname & \TextField[name=vorname, width=\linewidth, borderstyle=U]{} & \\[11pt]
  Nachname & \TextField[name=nachname, width=\linewidth, borderstyle=U]{} & \\[11pt]
  Geburtsdatum & \TextField[name=geburtsdatum, width=\linewidth, borderstyle=U]{} & \\[11pt]
  Staatsangehörigkeit & \TextField[name=staatsangehoerigkeit, width=\linewidth, borderstyle=U]{} & \\[4pt]
  \hline
\end{tblr}

\vspace{4pt}

% === KONTAKTINFORMATIONEN (bordered box) ===
\begin{tblr}{
  width = \textwidth,
  colspec = { |X| },
  rowsep = 0pt,
  colsep = 0pt,
}
  \hline
  \begin{tblr}{
    width = \linewidth,
    colspec = { X },
  }
    \textbf{Kontaktinformationen}
  \end{tblr}
  \\
  \begin{tblr}{
    width = \linewidth,
    colspec = { Q[l] X[c]},
  }
    Straße und Hausnummer: & \TextField[name=strasse, width=\linewidth, borderstyle=U]{} \\[10pt]
    Postleitzahl, Ort: & \TextField[name=plzOrt, width=\linewidth, borderstyle=U]{}
  \end{tblr}
  \\[10pt]
  \begin{tblr}{
    width = \linewidth,
    colspec = { Q[l] X[c,1] Q[l] X[l,1.5] },
    rowsep = 0pt,
    row{odd}  = {b, belowsep = -15pt},   % ← applies to 1,3,5 automatically
    row{even} = {t, abovesep = -15pt},   % ← applies to 2,4,6 automatically
  }
  Telefon: &
    \TextField[name=telefon, width=\linewidth, borderstyle=U]{} & & \\
  &
    {\scriptsize Angaben freiwillig} & & \\[10pt]
  Fax: &
    \TextField[name=fax, width=\linewidth, borderstyle=U]{}
  & E-Mail: &
    \TextField[name=email, width=\linewidth, borderstyle=U]{} \\
  &
    {\scriptsize Angaben freiwillig}
  & &
    {\centering\scriptsize Angaben freiwillig} \\[10pt]
  Handy: &
    \TextField[name=handy, width=\linewidth, borderstyle=U]{}
  & Internetzugang: & \ChoiceMenu[radio, name=internetzugang]{}{{}=ja} ja \quad \ChoiceMenu[radio, name=internetzugang]{}{{}=nein} nein \\
  &
    {\scriptsize Angaben freiwillig} & & \\
  \end{tblr}
  \\ \hline
\end{tblr}

\vspace{4pt}

% === KUNDENDATEN (bordered box) ===
\begin{tblr}{
  width = \textwidth,
  colspec = { |X| },
  rowsep = 0pt,
  colsep = 0pt,
}
  \hline
  \begin{tblr}{
    width = \linewidth,
    colspec = { X },
  }
    \textbf{Kundendaten}
  \end{tblr}
  \\
  \begin{tblr}{
    width = \linewidth,
    colspec = { Q[l] X[l] },
  }
    {
      Ich bin/werde arbeitslos\\
      seit/ab:
    }
    & \begin{minipage}[t]{6cm}
      \TextField[name=seit, width=\linewidth, borderstyle=U]{}
      \\[-4pt]
      \centering
      {\scriptsize Monat/Jahr}
    \end{minipage}
    \\
    {
      Letzte Tätigkeit/Ausbildung/\\
      Sonstiges von:
    }
    & \begin{minipage}[t]{6cm}
      \TextField[name=letzteTaetigkeit, width=\linewidth, borderstyle=U]{}
      \\[-4pt]
      \centering
      {\scriptsize Monat/Jahr}
    \end{minipage}
    \\
    als
    & \begin{minipage}[t]{\linewidth}
      \TextField[name=als, width=\linewidth, borderstyle=U]{}
      \\[-4pt]
      \centering
      {\scriptsize Tätigkeit, Name und Ort des Arbeitgeber/Schule/etc.}
    \end{minipage}
  \end{tblr}
  \\
  \begin{tblr}{
    width = \linewidth,
    colspec = { llX },
  }
    Beziehen Sie Arbeitslosengeld oder haben es beantragt?
    & \ChoiceMenu[radio, name=arbeitslosengeld]{}{{}=nein} nein
    & \ChoiceMenu[radio, name=arbeitslosengeld]{}{{}=ja} ja seit \TextField[name=arbeitslosengeldSeit, width=2.5cm, borderstyle=U]{}
  \end{tblr}
  \\ \hline
\end{tblr}

\vspace{4pt}

% === HEALTH / DISABILITY SECTION (bordered box) ===
\begin{tblr}{
  width = \textwidth,
  colspec = { |Xll| },
}
  \hline
  Ich habe gesundheitliche Einschränkungen
  & \ChoiceMenu[radio, name=gesundheitlicheEinschraenkungen]{}{{}=ja} ja
  & \ChoiceMenu[radio, name=gesundheitlicheEinschraenkungen]{}{{}=nein} nein
  \\
  Grad der Behinderung
  & \ChoiceMenu[radio, name=gradDerBehinderung]{}{{}=schwerbehindert} schwerbehindert
  & \ChoiceMenu[radio, name=gradDerBehinderung]{}{{}=gleichgestellt} gleichgestellt
  \\
  Ich habe einen Antrag auf berufliche Rehabilitation gestellt.
  & \ChoiceMenu[radio, name=antragBeruflicheRehabilitation]{}{{}=ja} ja
  & \ChoiceMenu[radio, name=antragBeruflicheRehabilitation]{}{{}=nein} nein
  \\
  Liegt ein Bescheid vor?
  & \ChoiceMenu[radio, name=bescheidVorhanden]{}{{}=ja} ja
  & \ChoiceMenu[radio, name=bescheidVorhanden]{}{{}=nein} nein
  \\ \hline
\end{tblr}

\newpage

% === TEIL 2 HEADING + SCHULABSCHLUSS (single bordered box) ===

\formsection{Teil 2 \textendash{} Berufliche Daten (bei Vorlage eines aktuellen und vollständigen
  Lebenslaufes brauchen Sie die Angaben zur beruflichen Aus- und Weiterbildung
  und zum beruflichen Werdegang nicht ausfüllen)}

\begin{tblr}{
  width = \textwidth,
  colspec = { |lX| },
}
  \hline
  \SetCell[c=2]{l, font=\bfseries}
    Angaben zu Ihrem Schulabschluss
    \\[8pt]
  \CheckBox[name=keinSchulabschluss]{}\ kein Schulabschluss & \CheckBox[name=fachhochschulreife]{}\ Fachhochschulreife
  \\[4pt]
  \CheckBox[name=abschlussFoerdershule]{}\ Abschluss Förderschule & \CheckBox[name=fachabitur]{}\ Fachabitur
  \\[4pt]
  \CheckBox[name=hauptschulabschluss]{}\ Hauptschulabschluss & \CheckBox[name=abitur]{}\ Abitur
  \\[4pt]
  \CheckBox[name=mittlereReife]{}\ Mittlere Reife
  \\[8pt]
  \SetCell[c=2]{l}
    \textbf{Schulentlassungsjahr} \quad \TextField[name=schulentlassungsjahr, width=3cm, borderstyle=U]{}
  \\ \hline
\end{tblr}

\vspace{8pt}

% === BERUFLICHE AUS- UND WEITERBILDUNG (bold heading) ===
\newcommand{\tablestrut}{\rule{0pt}{20pt}}
\newcommand{\rowheight}{4em}
\begin{tblr}{
  width = \textwidth,
  colspec = { |X| },
  rowsep = 0pt,
  colsep = 0pt,
}
  \hline
  \begin{tblr}{
    width = \linewidth,
    colspec = { X },
  }
    \textbf{Angaben zu Ihrer beruflichen Aus- und Weiterbildung}
  \end{tblr}
  \\
  \begin{tblr}{
    width = \linewidth,
    colspec = {
      Q[c,m,colsep=0pt]|
      Q[c,m,colsep=0pt]|
      X[c,m,colsep=0pt]|
      X[c,m,colsep=0pt]|
      Q[c, m, wd=2em]
      Q[c, m, wd=2em]
    },
    row{2} = {font=\scriptsize, valign=f},
    row{3-5} = {valign=m, rowsep=0pt},
  }
    \SetCell[c=2]{c} Zeitraum & & Ausbildungsstätte & Ausbildung als & \SetCell[c=2]{c} Abschluss\\
    \hline
    {
      von\\Tag/Monat/Jahr
    } &
    {
      bis\\Tag/Monat/Jahr
    } &
    (Institution bzw. Unternehmen, Ort) &
    Bezeichnung &
    ja &
    nein
    \\ \hline
    \begin{minipage}[c][\rowheight][c]{6em}
      \TextField[name=ausbildungOne, width=\linewidth, borderwidth=0]{}
    \end{minipage}
    & \begin{minipage}[c][\rowheight][c]{6em}
      \TextField[name=ausbildungOneBis, width=\linewidth, borderwidth=0]{}
    \end{minipage}
    & \begin{minipage}[c][\rowheight][c]{\linewidth}
      \TextField[name=ausbildungsstaetteOne, width=\linewidth, multiline=true, height=\rowheight, borderwidth=0]{}
    \end{minipage}
    & \begin{minipage}[c][\rowheight][c]{\linewidth}
      \TextField[name=ausbildungOneBezeichnung, width=\linewidth, multiline=true, height=\rowheight, borderwidth=0]{}
    \end{minipage}
    & \begin{minipage}[c][\rowheight][c]{2em}
      \centering
      \ChoiceMenu[radio, name=abschlussOne]{}{{}=ja}
    \end{minipage}
    & \begin{minipage}[c][\rowheight][c]{2em}
      \centering
      \ChoiceMenu[radio, name=abschlussOne]{}{{}=nein}
    \end{minipage}
    \\ \hline
    \begin{minipage}[c][\rowheight][c]{6em}
      \TextField[name=ausbildungTwo, width=\linewidth, borderwidth=0]{}
    \end{minipage}
    & \begin{minipage}[c][\rowheight][c]{6em}
      \TextField[name=ausbildungTwoBis, width=\linewidth, borderwidth=0]{}
    \end{minipage}
    & \begin{minipage}[c][\rowheight][c]{\linewidth}
      \TextField[name=ausbildungsstaetteTwo, width=\linewidth, multiline=true, height=\rowheight, borderwidth=0]{}
    \end{minipage}
    & \begin{minipage}[c][\rowheight][c]{\linewidth}
      \TextField[name=ausbildungTwoBezeichnung, width=\linewidth, multiline=true, height=\rowheight, borderwidth=0]{}
    \end{minipage}
    & \begin{minipage}[c][\rowheight][c]{2em}
      \centering
      \ChoiceMenu[radio, name=abschlussTwo]{}{{}=ja}
    \end{minipage}
    & \begin{minipage}[c][\rowheight][c]{2em}
      \centering
      \ChoiceMenu[radio, name=abschlussTwo]{}{{}=nein}
    \end{minipage}
    \\ \hline
    \begin{minipage}[c][\rowheight][c]{6em}
      \TextField[name=ausbildungThree, width=\linewidth, borderwidth=0]{}
    \end{minipage}
    & \begin{minipage}[c][\rowheight][c]{6em}
      \TextField[name=ausbildungThreeBis, width=\linewidth, borderwidth=0]{}
    \end{minipage}
    & \begin{minipage}[c][\rowheight][c]{\linewidth}
      \TextField[name=ausbildungsstaetteThree, width=\linewidth, multiline=true, height=\rowheight, borderwidth=0]{}
    \end{minipage}
    & \begin{minipage}[c][\rowheight][c]{\linewidth}
      \TextField[name=ausbildungThreeBezeichnung, width=\linewidth, multiline=true, height=\rowheight, borderwidth=0]{}
    \end{minipage}
    & \begin{minipage}[c][\rowheight][c]{2em}
      \centering
      \ChoiceMenu[radio, name=abschlussThree]{}{{}=ja}
    \end{minipage}
    & \begin{minipage}[c][\rowheight][c]{2em}
      \centering
      \ChoiceMenu[radio, name=abschlussThree]{}{{}=nein}
    \end{minipage}
    \\ \hline
  \end{tblr}
\end{tblr}

\vspace{8pt}

% === BERUFLICHER WERDEGANG (bold heading, no box) ===
\renewcommand{\rowheight}{2.5em}
\begin{tblr}{
  width = \textwidth,
  colspec = { |X| },
  rowsep = 0pt,
  colsep = 0pt,
}
  \hline
  \begin{tblr}{
    width = \linewidth,
    colspec = { X },
  }
    \textbf{Angaben zu Ihrem beruflichen Werdegang}
  \end{tblr}
  \\
  \begin{tblr}{
    width = \linewidth,
    colspec = {
      c|
      c|
      Q[c, co=1]|
      Q[c, co=1]
    },
    colsep = 0pt,
    row{2} = {font=\scriptsize, valign=f},
    row{3-8} = {valign=m, rowsep=0pt},
  }
    \SetCell[c=2]{c} Zeitraum & & Beschäftigungsstelle & Tätigkeit als
    \\ \hline
    {
      von\\Tag/Monat/Jahr
    } &
    {
      bis\\Tag/Monat/Jahr
    } &
    (Name, Ort) &
    Bezeichnung
    \\ \hline
    \begin{minipage}[c][\rowheight][c]{6em}
      \TextField[name=beschaeftigungOne, width=\linewidth, borderwidth=0]{}
    \end{minipage}
    & \begin{minipage}[c][\rowheight][c]{6em}
      \TextField[name=beschaeftigungOneBis, width=\linewidth, borderwidth=0]{}
    \end{minipage}
    & \begin{minipage}[c][\rowheight][c]{\linewidth}
      \TextField[name=beschaeftigungsstelleOne, width=\linewidth, multiline=true, height=\rowheight, borderwidth=0]{}
    \end{minipage}
    & \begin{minipage}[c][\rowheight][c]{\linewidth}
      \TextField[name=beschaeftigungOneTätigkeit, width=\linewidth, multiline=true, height=\rowheight, borderwidth=0]{}
    \end{minipage}
    \\ \hline
    \begin{minipage}[c][\rowheight][c]{6em}
      \TextField[name=beschaeftigungTwo, width=\linewidth, borderwidth=0]{}
    \end{minipage}
    & \begin{minipage}[c][\rowheight][c]{6em}
      \TextField[name=beschaeftigungTwoBis, width=\linewidth, borderwidth=0]{}
    \end{minipage}
    & \begin{minipage}[c][\rowheight][c]{\linewidth}
      \TextField[name=beschaeftigungsstelleTwo, width=\linewidth, multiline=true, height=\rowheight, borderwidth=0]{}
    \end{minipage}
    & \begin{minipage}[c][\rowheight][c]{\linewidth}
      \TextField[name=beschaeftigungTwoTätigkeit, width=\linewidth, multiline=true, height=\rowheight, borderwidth=0]{}
    \end{minipage}
    \\ \hline
    \begin{minipage}[c][\rowheight][c]{6em}
      \TextField[name=beschaeftigungThree, width=\linewidth, borderwidth=0]{}
    \end{minipage}
    & \begin{minipage}[c][\rowheight][c]{6em}
      \TextField[name=beschaeftigungThreeBis, width=\linewidth, borderwidth=0]{}
    \end{minipage}
    & \begin{minipage}[c][\rowheight][c]{\linewidth}
      \TextField[name=beschaeftigungsstelleThree, width=\linewidth, multiline=true, height=\rowheight, borderwidth=0]{}
    \end{minipage}
    & \begin{minipage}[c][\rowheight][c]{\linewidth}
      \TextField[name=beschaeftigungThreeTätigkeit, width=\linewidth, multiline=true, height=\rowheight, borderwidth=0]{}
    \end{minipage}
    \\ \hline
    \begin{minipage}[c][\rowheight][c]{6em}
      \TextField[name=beschaeftigungFour, width=\linewidth, borderwidth=0]{}
    \end{minipage}
    & \begin{minipage}[c][\rowheight][c]{6em}
      \TextField[name=beschaeftigungFourBis, width=\linewidth, borderwidth=0]{}
    \end{minipage}
    & \begin{minipage}[c][\rowheight][c]{\linewidth}
      \TextField[name=beschaeftigungsstelleFour, width=\linewidth, multiline=true, height=\rowheight, borderwidth=0]{}
    \end{minipage}
    & \begin{minipage}[c][\rowheight][c]{\linewidth}
      \TextField[name=beschaeftigungFourTätigkeit, width=\linewidth, multiline=true, height=\rowheight, borderwidth=0]{}
    \end{minipage}
    \\ \hline
    \begin{minipage}[c][\rowheight][c]{6em}
      \TextField[name=beschaeftigungFive, width=\linewidth, borderwidth=0]{}
    \end{minipage}
    & \begin{minipage}[c][\rowheight][c]{6em}
      \TextField[name=beschaeftigungFiveBis, width=\linewidth, borderwidth=0]{}
    \end{minipage}
    & \begin{minipage}[c][\rowheight][c]{\linewidth}
      \TextField[name=beschaeftigungsstelleFive, width=\linewidth, multiline=true, height=\rowheight, borderwidth=0]{}
    \end{minipage}
    & \begin{minipage}[c][\rowheight][c]{\linewidth}
      \TextField[name=beschaeftigungFiveTätigkeit, width=\linewidth, multiline=true, height=\rowheight, borderwidth=0]{}
    \end{minipage}
    \\ \hline
    \begin{minipage}[c][\rowheight][c]{6em}
      \TextField[name=beschaeftigungSix, width=\linewidth, borderwidth=0]{}
    \end{minipage}
    & \begin{minipage}[c][\rowheight][c]{6em}
      \TextField[name=beschaeftigungSixBis, width=\linewidth, borderwidth=0]{}
    \end{minipage}
    & \begin{minipage}[c][\rowheight][c]{\linewidth}
      \TextField[name=beschaeftigungsstelleSix, width=\linewidth, multiline=true, height=\rowheight, borderwidth=0]{}
    \end{minipage}
    & \begin{minipage}[c][\rowheight][c]{\linewidth}
      \TextField[name=beschaeftigungSixTätigkeit, width=\linewidth, multiline=true, height=\rowheight, borderwidth=0]{}
    \end{minipage}
    \\ \hline
  \end{tblr}
\end{tblr}

\vspace{8pt}

% === MOBILITäT (bordered box with heading inside) ===
  \begin{tblr}{
    width = \textwidth,
    colspec = { |l c c X| },
  }
    \hline
    \SetCell[c=4]{l, font=\bfseries}
      Angaben zu Ihrer Mobilität
      & & & \\
    \hline
    Führerschein und Fahrzeug & Führerschein & Fahrzeug vorhanden \\
    Kraftrad\tablestrut & \CheckBox[name=kraftradFuehrerschein]{} & \CheckBox[name=kraftradFahrzeug]{} \\
    PKW\tablestrut & \CheckBox[name=pkwFuehrerschein]{} & \CheckBox[name=pkwFahrzeug]{} \\
    LKW\tablestrut & \CheckBox[name=lkwFuehrerschein]{} & \\
    \hline
  \end{tblr}

\newpage

% === SPRACHKENNTNISSEN (bordered box, at top of page 3) ===
  \begin{tblr}{
    width = \textwidth,
    colspec = { |l c c X| }
  }
    \hline
    \SetCell[c=4]{l, font=\bfseries}
      Angaben zu Ihren Sprachkenntnissen
      & & & \\
    \hline
    Muttersprache\tablestrut & \SetCell[c=3]{l} \TextField[name=muttersprache, width=8cm, borderstyle=U]{} & & \\
    & Grundkenntnisse & gut & \\
    Deutsch & \ChoiceMenu[radio, name=deutsch]{}{{}=grundkenntnisse} & \ChoiceMenu[radio, name=deutsch]{}{{}=gut} & \\
    Englisch & \ChoiceMenu[radio, name=englisch]{}{{}=grundkenntnisse} & \ChoiceMenu[radio, name=englisch]{}{{}=gut} & \\
    \hline
  \end{tblr}

\vspace{12pt}

% === TEIL 3 HEADING (bold text, NOT boxed, en-dash) ===
\formsection{Teil 3 \textendash{} Vorbereitung Vermittlungsgespräch}

% === LARGE BORDERED BOX containing Bewerbung + Familiäres Umfeld ===
\begin{tblr}{
  width = \textwidth,
  colspec = { |X| },
  rowsep = 0pt,
  colsep = 0pt,
}
  \hline
  \begin{tblr}{
    width = \linewidth,
    colspec = { lX },
  }
    \tablestrut \textbf{Ich suche eine Stelle als}
    & \TextField[name=stelle, width=\linewidth, borderstyle=U]{}
  \end{tblr}
  \\ \hline
  \begin{tblr}{
    width = \linewidth,
    colspec = {
      X[colsep=0pt]|
      X[colsep=0pt]|
      c|
      c
    },
    vline{3} = {1}{0pt},
    rows = {valign = m},
    row{3-8} = {rowsep = 0pt},
  }
    \SetCell[c=2]{l} \hspace{6pt}\textbf{Ich habe mich bisher beworben} \tablestrut & & \SetCell[c=2]{c} \textbf{Ergebnis} \\
    \hspace{6pt}bei & \hspace{6pt}als & offen & Absage \\ \hline
    \begin{minipage}[c][\rowheight][c]{\linewidth}
      \TextField[name=bewerbungOne, width=\linewidth, multiline=true, height=\rowheight, borderwidth=0]{}
    \end{minipage}
    & \begin{minipage}[c][\rowheight][c]{\linewidth}
      \TextField[name=bewerbungOneTaetigkeit, width=\linewidth, multiline=true, height=\rowheight, borderwidth=0]{}
    \end{minipage}
    & \ChoiceMenu[radio, name=bewerbungOffenOne]{}{{}=offen}
    & \ChoiceMenu[radio, name=bewerbungOffenOne]{}{{}=absage}
    \\ \hline
    \begin{minipage}[c][\rowheight][c]{\linewidth}
      \TextField[name=bewerbungTwo, width=\linewidth, multiline=true, height=\rowheight, borderwidth=0]{}
    \end{minipage}
    & \begin{minipage}[c][\rowheight][c]{\linewidth}
      \TextField[name=bewerbungTwoTaetigkeit, width=\linewidth, multiline=true, height=\rowheight, borderwidth=0]{}
    \end{minipage}
    & \ChoiceMenu[radio, name=bewerbungOffenTwo]{}{{}=offen}
    & \ChoiceMenu[radio, name=bewerbungOffenTwo]{}{{}=absage}
    \\ \hline
    \begin{minipage}[c][\rowheight][c]{\linewidth}
      \TextField[name=bewerbungThree, width=\linewidth, multiline=true, height=\rowheight, borderwidth=0]{}
    \end{minipage}
    & \begin{minipage}[c][\rowheight][c]{\linewidth}
      \TextField[name=bewerbungThreeTaetigkeit, width=\linewidth, multiline=true, height=\rowheight, borderwidth=0]{}
    \end{minipage}
    & \ChoiceMenu[radio, name=bewerbungOffenThree]{}{{}=offen}
    & \ChoiceMenu[radio, name=bewerbungOffenThree]{}{{}=absage}
    \\ \hline
    \begin{minipage}[c][\rowheight][c]{\linewidth}
      \TextField[name=bewerbungFour, width=\linewidth, multiline=true, height=\rowheight, borderwidth=0]{}
    \end{minipage}
    & \begin{minipage}[c][\rowheight][c]{\linewidth}
      \TextField[name=bewerbungFourTaetigkeit, width=\linewidth, multiline=true, height=\rowheight, borderwidth=0]{}
    \end{minipage}
    & \ChoiceMenu[radio, name=bewerbungOffenFour]{}{{}=offen}
    & \ChoiceMenu[radio, name=bewerbungOffenFour]{}{{}=absage}
    \\ \hline
    \begin{minipage}[c][\rowheight][c]{\linewidth}
      \TextField[name=bewerbungFive, width=\linewidth, multiline=true, height=\rowheight, borderwidth=0]{}
    \end{minipage}
    & \begin{minipage}[c][\rowheight][c]{\linewidth}
      \TextField[name=bewerbungFiveTaetigkeit, width=\linewidth, multiline=true, height=\rowheight, borderwidth=0]{}
    \end{minipage}
    & \ChoiceMenu[radio, name=bewerbungOffenFive]{}{{}=offen}
    & \ChoiceMenu[radio, name=bewerbungOffenFive]{}{{}=absage}
    \\ \hline
    \begin{minipage}[c][\rowheight][c]{\linewidth}
      \TextField[name=bewerbungSix, width=\linewidth, multiline=true, height=\rowheight, borderwidth=0]{}
    \end{minipage}
    & \begin{minipage}[c][\rowheight][c]{\linewidth}
      \TextField[name=bewerbungSixTaetigkeit, width=\linewidth, multiline=true, height=\rowheight, borderwidth=0]{}
    \end{minipage}
    & \ChoiceMenu[radio, name=bewerbungOffenSix]{}{{}=offen}
    & \ChoiceMenu[radio, name=bewerbungOffenSix]{}{{}=absage}
  \end{tblr}
  \\ \hline
  \begin{tblr}{
    width = \linewidth,
    colspec = { lX },
  }
    \tablestrut Bitte legen Sie meine Bewerbung nicht vor bei (Firma):
    & \TextField[name=firma, width=\linewidth, borderstyle=U]{}
  \end{tblr}
  \\ \hline
  \begin{tblr}{
    width = \linewidth,
    colspec = { X },
  }
    \textbf{Familiäres Umfeld} \tablestrut
  \end{tblr}
  \\ \hline
  \begin{tblr}{
    width = \linewidth,
    colspec = { lllX },
  }
    \SetCell[c=4]{l}
      \textbf{In meinem familiären Umfeld sind folgende Gegebenheiten zu berücksichtigen:} \tablestrut
      & & & &
    \\
    & & Betreuung sichergestellt & \\
    \CheckBox[name=kinderUnter15Jahren]{} & Kinder unter 15 Jahren im Haushalt
    & \ChoiceMenu[radio, name=betreuungSichergestellt]{}{{}=ja} ja \hfill \ChoiceMenu[radio, name=betreuungSichergestellt]{}{{}=nein} nein \\
    & Anzahl \TextField[name=anzahlKinder, width=3cm, borderstyle=U]{} \\[6pt]
    \CheckBox[name=pflegeAngehoeriger]{} & Pflege eines Angehörigen & & \\[6pt]
    \CheckBox[name=sonstigeVerpflichtungen]{} & Sonstige Verpflichtungen & &
  \end{tblr}
  \\ \hline
\end{tblr}

\newpage

\formsection{Teil 4 \textendash{} Angaben zu Ihrem Migrationshintergrund}

\begin{tblr}{
  width = \textwidth,
  colspec = { |llX[l]ccc| },
  rows = {rowsep = 10pt, valign = m},
}
  \hline
  & & & Ja & Nein & {Keine\\Angabe}
  \\ \hline
  1 & \SetCell[c=2]{l} Ich besitze die deutsche Staatsangehörigkeit. &
  & \ChoiceMenu[radio, name=q1]{}{{}=ja}
  & \ChoiceMenu[radio, name=q1]{}{{}=nein}
  & \ChoiceMenu[radio, name=q1]{}{{}=keine}
  \\ \cline{2-6}
  \SetCell{valign = h}2 & \SetCell{valign = h}a) &
  {
    Ich bin in Deutschland geboren.\\
    (heutiges Gebiet der Bundesrepublik Deutschland)
  }
  & \ChoiceMenu[radio, name=q2a]{}{{}=ja}
  & \ChoiceMenu[radio, name=q2a]{}{{}=nein}
  & \ChoiceMenu[radio, name=q2a]{}{{}=keine}
  \\
  & \SetCell{valign = h}b) &
  {
    Ich bin nach 1949 nach Deutschland zugewandert.\\
    (heutiges Gebiet der Bundesrepublik Deutschland)
  }
  & \ChoiceMenu[radio, name=q2b]{}{{}=ja}
  & \ChoiceMenu[radio, name=q2b]{}{{}=nein}
  & \ChoiceMenu[radio, name=q2b]{}{{}=keine}
  \\ \cline{2-6}
  \SetCell{valign = h}3
  & \SetCell{valign = h}a)
  & Mein Vater ist außerhalb Deutschlands (heutiges Gebiet der Bundesrepublik Deutschland) geboren und nach 1949 zugewandert.
  & \ChoiceMenu[radio, name=q3a]{}{{}=ja}
  & \ChoiceMenu[radio, name=q3a]{}{{}=nein}
  & \ChoiceMenu[radio, name=q3a]{}{{}=keine}
  \\
  & \SetCell{valign = h}b)
  & Meine Mutter ist außerhalb Deutschlands (heutiges Gebiet der Bun\-des\-re\-pu\-blik Deutschland) geboren und nach 1949 zugewandert.
  & \ChoiceMenu[radio, name=q3b]{}{{}=ja}
  & \ChoiceMenu[radio, name=q3b]{}{{}=nein}
  & \ChoiceMenu[radio, name=q3b]{}{{}=keine}
  \\ \hline
\end{tblr}

\vspace{40pt}

  \begin{tblr}{
    width = \textwidth,
    colspec = { |llXccc| },
    rows = {rowsep = 10pt, valign = m},
  }
    \hline
    \SetCell[c=6]{l}
    {
      Zusätzliche Fragen für deutsche Staatsangehörige, die nach 1949 zugewandert sind
      \\
      \textbf{(Frage 1 = \glqq Ja\grqq\ \underline{und} Frage 2b = \glqq Ja\grqq)}
    }
    \\ \hline
    & & & Ja & Nein & {Keine\\Angabe}
    \\ \hline
      \SetCell{valign = h}4 & \SetCell{valign = h}a)
      & Ich habe die deutsche Staatsangehörigkeit als Aussiedler oder Spätaussiedler erworben.
      & \ChoiceMenu[radio, name=q4a]{}{{}=ja}
      & \ChoiceMenu[radio, name=q4a]{}{{}=nein}
      & \ChoiceMenu[radio, name=q4a]{}{{}=keine}
    \\
      & \SetCell{valign = h}b)
      & Ich habe die deutsche Staatsangehörigkeit als Ehegatte / Ehegattin eines (Spät-)Aus\-sied\-lers erworben.
      & \ChoiceMenu[radio, name=q4b]{}{{}=ja}
      & \ChoiceMenu[radio, name=q4b]{}{{}=nein}
      & \ChoiceMenu[radio, name=q4b]{}{{}=keine}
    \\
      & \SetCell{valign = h}c)
      & Ich habe die deutsche Staatsangehörigkeit als Kind oder Enkelkind eines (Spät-)Aus\-sied\-lers erworben.
      & \ChoiceMenu[radio, name=q4c]{}{{}=ja}
      & \ChoiceMenu[radio, name=q4c]{}{{}=nein}
      & \ChoiceMenu[radio, name=q4c]{}{{}=keine}
    \\ \hline
  \end{tblr}%
\end{Form}

\end{document}
